\section{Fondamenti di Geometria Proiettiva}
\label{sec:geom_proj}

\subsection{Coordinate omogenee e spazio proiettivo}

Lo spazio proiettivo \(\mathbb{P}^n\) estende \(\mathbb{R}^n\) aggiungendo
punti all'infinito: un punto di \(\mathbb{P}^n\) è una classe di equivalenza
\([x_0 : x_1 : \cdots : x_n]\) con \((x_0,\dots,x_n)\neq\mathbf{0}\),
dove due vettori sono equivalenti se proporzionali.

In \(\mathbb{P}^2\) un punto euclideo \((x,y)\) si rappresenta come
\(\homo{x} = (x_1, x_2, x_3)^T\) con \(x = x_1/x_3\),
\(y = x_2/x_3\). Una retta di equazione \(ax+by+c=0\) corrisponde al
vettore \(\mathbf{l}=(a,b,c)^T\), e il punto \(\homo{x}\) giace su
\(\mathbf{l}\) se e solo se \(\mathbf{l}^T\homo{x} = 0\).

\begin{defn}[Dualità in \(\mathbb{P}^2\)]
In \(\mathbb{P}^2\) punti e rette sono \emph{duali}: ogni enunciato che
coinvolge punti e rette rimane valido scambiando i due ruoli.
L'intersezione di due rette \(\mathbf{l},\mathbf{m}\) è
\(\homo{x} = \mathbf{l} \times \mathbf{m}\);
la retta per due punti \(\homo{x},\homo{y}\) è
\(\mathbf{l} = \homo{x} \times \homo{y}\).
\end{defn}

In \(\mathbb{P}^3\) un punto è \(\homo{X} = (X,Y,Z,W)^T\); un piano è
il vettore duale \(\boldsymbol{\pi} = (\alpha,\beta,\gamma,\delta)^T\)
tale che \(\boldsymbol{\pi}^T \homo{X} = 0\)~\cite[Cap.~2]{HZ2004}.

\subsection{Punti propri e punti impropri}
\label{sub:punti_propri_impropri}

In \(\mathbb{P}^2\) un punto \(\homo{x} = (x_1, x_2, x_3)^T\) si dice
\emph{proprio} se \(x_3 \neq 0\): in tal caso la divisione per \(x_3\)
restituisce le coordinate euclidee finite \((x_1/x_3,\, x_2/x_3)\).
Se invece \(x_3 = 0\), il punto è \emph{improprio} (o \emph{punto
all'infinito}) e rappresenta una \emph{direzione} nel piano:
non ha controparte euclidea, ma è pienamente definito come
elementodi \(\mathbb{P}^2\).

\begin{oss}
Il piano all'infinito \(\pi_\infty\) di \(\mathbb{P}^3\) ha equazione
omogenea \(W = 0\) ed è il luogo dei punti impropri dello spazio;
in \(\mathbb{P}^2\) la retta all'infinito \(\mathbf{l}_\infty = (0,0,1)^T\)
è il luogo dei punti impropri del piano.
Rette parallele di direzione \(\mathbf{d} = (d_1, d_2)^T\) si incontrano
tutte nel punto improprio \((d_1, d_2, 0)^T \in \mathbf{l}_\infty\).
\end{oss}

Questa distinzione è operativamente rilevante per il progetto:
quando la camera osserva il pattern planare, la proiezione prospettica
trasforma punti propri del piano del target in punti propri
dell'immagine. Tuttavia la retta all'infinito \(\mathbf{l}_\infty\)
del piano del pattern viene mappata in una retta propria
dell'immagine (la \emph{retta orizzonte}), e i punti impropri del
pattern corrispondono ai \emph{punti di fuga} nell'immagine~\cite[Sez.~2.2]{HZ2004}.
Riportare l'immagine a una vista rettificata equivale quindi
a rimandare la retta orizzonte verso l'infinito, ripristinando
la struttura affine del piano originale.

\subsection{Trasformazioni proiettive e gerarchia}
\label{sub:proiettivita}

Una \emph{proiettività} (omografia) di \(\mathbb{P}^n\) è una biiezione
indotta da una matrice invertibile \((n{+}1)\times(n{+}1)\):
\[
  \homo{x}' = \mat{H}\,\homo{x}, \qquad \mat{H} \in GL(n+1).
\]

Poiché \(\mat{H}\) è invertibile, anche la trasformazione inversa
\(\homo{x} = \mat{H}^{-1}\homo{x}'\) è una proiettività. Questo fatto,
apparentemente ovvio, è geometricamente rilevante: l'omografia che
correla il piano del pattern con il piano immagine può essere
invertita per ottenere la vista rettificata, e tale inversione
è ancora una proiettività — nessuna nuova classe di trasformazione
è necessaria per descrivere la rettifica.

\begin{table}[H]
\centering
\caption{Gerarchia delle trasformazioni planari con gradi di libertà (dof)
         e invarianti geometrici \cite[Tab.~2.1]{HZ2004}.}
\label{tab:trasformazioni}
\begin{tabular}{llcc}
\toprule
\textbf{Tipo} & \textbf{Forma di \(\mat{H}\)} & \textbf{dof} & \textbf{Invarianti} \\
\midrule
Isometria    & \(\bigl(\begin{smallmatrix}\mat{R}&\mathbf{t}\\\mathbf{0}^T&1\end{smallmatrix}\bigr)\) & 3 & lunghezze, angoli \\
Similarità   & \(\bigl(\begin{smallmatrix}s\mat{R}&\mathbf{t}\\\mathbf{0}^T&1\end{smallmatrix}\bigr)\) & 4 & angoli, rapporti di lunghezza \\
Affinità     & \(\bigl(\begin{smallmatrix}\mat{A}&\mathbf{t}\\\mathbf{0}^T&1\end{smallmatrix}\bigr)\) & 6 & parallelismo, rapporti di aree \\
Proiettività & \(\mat{H}\) (generica) & 8 & birapporto, collinearità \\
\bottomrule
\end{tabular}
\end{table}

\subsection{Le affinità come caso speciale delle proiettività}
\label{sub:affinita_proiettivita}

Dall'ispezione della Tabella~\ref{tab:trasformazioni} emerge la gerarchia:
\[
  \text{Isometrie} \subset \text{Similarità} \subset \text{Affinità}
  \subset \text{Proiettività}.
\]
Una trasformazione affine ha la forma
\[
  \mat{H}_{\text{aff}} =
  \begin{pmatrix} \mat{A} & \mathbf{t} \\ \mathbf{0}^T & 1 \end{pmatrix},
  \qquad \mat{A} \in GL(2),
\]
dove l'ultima riga è fissa a \((0,0,1)\). Questo vincolo ha una
interpretazione geometrica precisa: una trasformazione affine
\emph{manda i punti impropri in punti impropri}, cioè preserva
la retta all'infinito \(\mathbf{l}_\infty\) come insieme
(non necessariamente punto per punto). In altre parole, le rette
parallele prima della trasformazione rimangono parallele dopo,
perché i loro punti di fuga restano su \(\mathbf{l}_\infty\).

Una proiettività generica, invece, può mappare punti impropri
in punti propri e viceversa: è esattamente questo che accade
quando la camera inquadra il pattern da una posa obliqua,
portando i punti di fuga del pattern all'interno dell'immagine.
La rettifica dell'immagine consiste nel trovare una omografia
\(\mat{H}\) che riporti quei punti di fuga all'infinito,
restaurando la struttura affine del piano originale.

In CCalibGA questa comprensione è fondamentale: la matrice
intrinseca \(\mat{K}\) e la posa \([\mat{R}|\mathbf{t}]\)
stimate durante la calibrazione permettono di costruire esattamente
quell'omografia di rettifica.

\subsection{Birapporto}

Il \emph{birapporto} (cross-ratio) di quattro punti collineari
\(A,B,C,D\) è l'unico invariante scalare delle proiettività:
\[
  (A,B;C,D)
  \;=\;
  \frac{\overline{AC}\cdot\overline{BD}}{\overline{AD}\cdot\overline{BC}},
\]
dove \(\overline{PQ}\) denota la distanza orientata tra \(P\) e \(Q\).
Il birapporto è preservato da qualsiasi proiettività e rappresenta
lo strumento canonico per verificare la coerenza delle corrispondenze
punto-punto nel matching tra immagini~\cite[Sez.~2.4]{HZ2004}.

\subsection{Conica assoluta e IAC}
\label{sub:IAC}

La \emph{conica assoluta} \(\Omega_\infty\) è definita come
l'intersezione del piano all'infinito \(\pi_\infty\) (con \(W=0\) in
\(\mathbb{P}^3\)) con la sfera unitaria; la sua matrice è
\(\mat{I}_3\)~\cite[Sez.~8.3]{HZ2004}. I suoi punti soddisfano:
\[
  X^2+Y^2+Z^2=0, \quad W=0,
\]
e sono necessariamente complessi; i \emph{punti circolari all'infinito}
\(\mathbf{I} = (1,i,0)^T\) e \(\mathbf{J}=(1,-i,0)^T\) ne sono la
proiezione nel piano.

L'\emph{Immagine della Conica Assoluta} (IAC), indicata con \(\omega\),
è la proiezione di \(\Omega_\infty\) tramite \(\mat{P}=\mat{K}[\mat{R}|\mathbf{t}]\):
\[
  \omega = \mat{K}^{-T}\mat{K}^{-1} = (\mat{K}\mat{K}^T)^{-1}.
  \label{eq:IAC}
\]
Questa relazione è il cuore del metodo di Zhang: imporre che
\(\mat{H}\) mappi \(\Omega_\infty\) in \(\omega\) fornisce due equazioni
lineari nei coefficienti di \(\omega\), da cui si estrae \(\mat{K}\) per
decomposizione di Cholesky~\cite[Prop.~8.5]{HZ2004}.

\begin{figure}[H]
\centering
\begin{tikzpicture}[scale=1.0, >=Stealth,
  plane/.style={draw=gray!60, fill=gray!10, opacity=0.8},
  arr/.style={->, thick, blue!70!black}]

  \draw[plane] (-1.5,-0.8) -- (1.5,-0.8) -- (2.5,0.5) -- (-0.5,0.5) -- cycle;
  \node[below right, gray] at (2.3,0.5) {\small Piano \(Z=0\)};

  \draw[plane, fill=orange!15]
    (3.5,-0.8) -- (5.5,-0.8) -- (5.5,1.6) -- (3.5,1.6) -- cycle;
  \node[above, orange!70!black] at (4.5,1.65) {\small Piano immagine};

  \node[circle, fill=black, inner sep=2pt,
        label=above:{\small\(\mat{C}\)}] (C) at (2.0,1.8) {};

  \foreach \px/\py/\qx/\qy in
    {0.0/-0.2/3.9/0.2, 0.5/0.3/4.3/0.9, -0.5/0.1/4.1/0.5}{
    \draw[arr] (C) -- (\px,\py);
    \draw[arr,dashed] (\px,\py) -- (\qx,\qy);
  }

  \node[below] at (0.0,-0.25) {\small \(\homo{X}\)};
  \node[right]  at (4.15,0.5) {\small \(\homo{x}\)};

  \draw[->, thick, red!70!black] (1.0,0.3) to[bend left=20]
    node[above, red!70!black] {\(\mat{H}\)} (3.9,0.5);
\end{tikzpicture}
\caption{Proiettività \(\mat{H}\colon\mathbb{P}^2\to\mathbb{P}^2\) indotta
         dalla camera: mette in corrispondenza il piano \(Z=0\) del target
         con il piano immagine.}
\label{fig:omografia}
\end{figure}
